\documentclass{article}
\usepackage{amsmath} % Para usar entornos matemáticos
\usepackage{amsfonts} % Para usar fuentes y símbolos matemáticos
\usepackage{amssymb} % Para usar símbolos matemáticos adicionales

\title{Análisis Tensorial: ecuaciones de Gauss-Weingarten}
\author{David E. Duarte}
\date{\today}

\begin{document}
	\maketitle
	
	En el análisis tensorial, las ecuaciones de Gauss-Weingarten son un conjunto de ecuaciones que relacionan las curvaturas intrínsecas de una superficie diferenciable en un espacio euclidiano tridimensional. Estas ecuaciones son fundamentales en la geometría diferencial y proporcionan información sobre la geometría local de la superficie.
	
	Las ecuaciones de Gauss-Weingarten se pueden expresar de la siguiente manera:
	
	\begin{enumerate}
		\item Ecuación de Gauss: $K = \text{det}(S)$
		\item Ecuación de Codazzi: $\nabla_X S - \nabla_{\nabla_X Y} N = 0$
		\item Ecuación de la curvatura media: $H = \frac{1}{2} \text{tr}(S)$
	\end{enumerate}
	
	Donde:
	\begin{itemize}
		\item $K$ es la curvatura gaussiana,
		\item $S$ es el tensor de forma,
		\item $N$ es el vector normal unitario a la superficie,
		\item $H$ es la curvatura media,
		\item $\nabla$ es el operador de derivada covariante en la superficie,
		\item $X$ y $Y$ son campos vectoriales tangentes a la superficie.
	\end{itemize}
	
	Estas ecuaciones son utilizadas para estudiar la geometría de superficies y proporcionan información sobre la curvatura de la superficie en cada punto.
	
	Un ejemplo práctico de aplicación de las ecuaciones de Gauss-Weingarten es en la teoría de las superficies mínimas. En este contexto, se busca encontrar superficies que tienen curvatura media igual a cero en cada punto, lo que implica que la ecuación de la curvatura media se reduce a $H = 0$. Al resolver las ecuaciones de Gauss-Weingarten para este caso particular, se pueden obtener condiciones necesarias para determinar si una superficie es mínima.
	
	Otro ejemplo práctico es en la teoría de la elasticidad, donde las ecuaciones de Gauss-Weingarten se utilizan para modelar la deformación de materiales sólidos. Al considerar una superficie que representa la frontera de un sólido deformable, las ecuaciones de Gauss-Weingarten pueden proporcionar información sobre las tensiones y deformaciones en la superficie.
	
	Las ecuaciones de Gauss-Weingarten son herramientas poderosas en el análisis tensorial que permiten estudiar la geometría de superficies y su comportamiento en diferentes contextos, desde la geometría diferencial hasta la física de materiales.
	
	\section{Caso con una spuerficie esférica}
	
	Para una superficie esférica, podemos aplicar las ecuaciones de Gauss-Weingarten para calcular la curvatura de la superficie en cada punto. En el caso de una superficie parametrizada como una esfera, las ecuaciones de Gauss-Weingarten se simplifican debido a la forma específica de la parametrización.
	
	Las ecuaciones de Gauss-Weingarten son las siguientes:
	
	1. \textbf{Ecuación de Gauss}:
	$$ K = \frac{LN - M^2}{EG - F^2} $$
	
	2. \textbf{Ecuación de la curvatura media}:
	$$ H = \frac{1}{2}(H_1 + H_2) = \frac{1}{2}(L + E) $$
	
	Para la esfera parametrizada que hemos considerado, podemos calcular los coeficientes de la primera forma fundamental $ E, F, G $ y los coeficientes de la segunda forma fundamental $ L, M, N $ en cada punto utilizando las derivadas parciales de la parametrización. Luego, podemos aplicar las ecuaciones de Gauss-Weingarten para obtener la curvatura gaussiana $ K $ y la curvatura media $ H $ en cada punto de la superficie esférica.
	
	En el caso de la esfera, la curvatura gaussiana $ K $ es constante y es igual a $ \frac{1}{R^2} $ (donde $ R $ es el radio de la esfera). La curvatura media $ H $ es constante y es igual a $ \frac{1}{R} $.
	
	Por lo tanto, al aplicar las ecuaciones de Gauss-Weingarten en el ejemplo de la esfera, podemos calcular la curvatura gaussiana y la curvatura media en cada punto de la superficie esférica, lo que nos proporciona información sobre la geometría de la superficie en términos de curvatura.
	
	\section{Ejemplo:}
	
	Vamos a considerar una esfera para este caso.
	
	1. \textbf{Superficie de una esfera}:
	La ecuación de una esfera en coordenadas cartesianas es:
	$$ x^2 + y^2 + z^2 = R^2 $$
	donde $ R $ es el radio de la esfera.
	
	2. \textbf{Parametrización de la esfera}:
	Podemos parametrizar la esfera con coordenadas esféricas $ (r, \theta, \phi) $ de la siguiente manera:
	$$ x = R \sin(\theta) \cos(\phi) $$
	$$ y = R \sin(\theta) \sin(\phi) $$
	$$ z = R \cos(\theta) $$
	donde $ 0 \leq \theta \leq \pi $ y $ 0 \leq \phi \leq 2\pi $.
	
	Veamos el ejemplo para una esfera de radio $ R = 3 $. En este caso, la ecuación de la esfera sería:
	$$ x^2 + y^2 + z^2 = 3^2 $$
	$$ x^2 + y^2 + z^2 = 9 $$
	
	La parametrización de la esfera con radio $ R = 3 $ en coordenadas esféricas sería:
	$$ x = 3 \sin(\theta) \cos(\phi) $$
	$$ y = 3 \sin(\theta) \sin(\phi) $$
	$$ z = 3 \cos(\theta) $$
	donde $ 0 \leq \theta \leq \pi $ y $ 0 \leq \phi \leq 2\pi $.
	
	La curvatura gaussiana $ K $ para una esfera de radio $ R $ está dada por:
	$$ K = \frac{1}{R^2} $$
	En este caso, para una esfera de radio 3, la curvatura gaussiana sería:
	$$ K = \frac{1}{3^2} = \frac{1}{9} $$
	
	La curvatura media $ H $ para una esfera de radio $ R $ está dada por:
	$$ H = \frac{1}{R} $$
	Para una esfera de radio 3, la curvatura media sería:
	$$ H = \frac{1}{3}  $$
	
\subsection{Ejemplo de codigo en Python que calcula los parámetros de curvatura para una Superficie Esférica}

\begin{verbatim}
	import sympy as sp
	
	# Ingresar  el radio de la esfera
	R = float(input("Ingrese el radio de la esfera: "))
	
	#  parametrizaciones para una esfera en coordenadas esféricas
	theta = sp.Symbol('theta')
	phi = sp.Symbol('phi')
	
	#  derivadas parciales de las parametrizaciones para una esfera de radio R
	dx_dtheta = R * sp.cos(theta) * sp.cos(phi)
	dx_dphi = -R * sp.sin(theta) * sp.sin(phi)
	dy_dtheta = R * sp.cos(theta) * sp.sin(phi)
	dy_dphi = R * sp.sin(theta) * sp.cos(phi)
	dz_dtheta = -R * sp.sin(theta)
	dz_dphi = 0
	
	#  coeficientes de la primera forma fundamental
	E = sp.simplify(sp.simplify(dx_dtheta**2 + dy_dtheta**2 + dz_dtheta**2))
	F = sp.simplify(sp.simplify(dx_dtheta*dx_dphi + dy_dtheta*dy_dphi + dz_dtheta*dz_dphi))
	G = sp.simplify(sp.simplify(dx_dphi**2 + dy_dphi**2 + dz_dphi**2))
	
	#  coeficientes de la segunda forma fundamental
	L = sp.simplify(sp.diff(dx_dphi, theta) + sp.diff(dy_dphi, theta) + sp.diff(dz_dphi, theta))
	M = sp.simplify(sp.diff(dx_dphi, phi) + sp.diff(dy_dphi, phi) + sp.diff(dz_dphi, phi))
	N = sp.simplify(sp.diff(dx_dtheta, phi) + sp.diff(dy_dtheta, phi) + sp.diff(dz_dtheta, phi))
	
	#  curvatura gaussiana y la curvatura media
	K = sp.simplify((L*N - M**2) / (E*G - F**2))
	H = sp.simplify((E*N + G*L - 2*F*M) / (2*(E*G - F**2)))
	
	print("\nCoeficientes de la primera forma fundamental:")
	print("E =", E)
	print("F =", F)
	print("G =", G)
	print("\nCoeficientes de la segunda forma fundamental:")
	print("L =", L)
	print("M =", M)
	print("N =", N)
	print("\nCurvatura gaussiana (K):")
	print("K =", K)
	print("\nCurvatura media (H):")
	print("H =", H)
\end{verbatim}
	
	
	
	
	
\end{document}